\section{完全正则空间上的测度}

若 \(\mathbf{T}\) 是拓扑空间，且 \(\overline{\mathbf{F}}\) 是 Banach 空间，则记号 \(\mathcal{C}^b (\mathrm{T};\mathbf{F})\) 表示定义在 \(\mathbf{T}\) 上、取值于 \(\mathbf{F}\) 的有界连续函数空间，并赋予一致收敛范数。当 \(\mathbf{F} = \mathbf{R}\) 时，此记号简写为 \(\mathcal{C}^b (\mathrm{T})\)；若无歧义，也简写为 \(\mathcal{C}^b\)，并以 \(\mathcal{C}_+^b (\mathrm{T})\) 或 \(\mathcal{C}_+^b\) 表示 \(\mathcal{C}^b (\mathrm{T})\) 中正函数所成的锥。\(\mathbf{T}\) 上有界复测度空间记为 \(\mathcal{M}^b (\mathrm{T};\mathbf{C})\)，有界实测度空间记为 \(\mathcal{M}^b (\mathrm{T})\) 或 \(\mathcal{M}^b\)，有界正测度所成的锥记为 \(\mathcal{M}_+^b (\mathrm{T})\) 或 \(\mathcal{M}_+^b\)。

